\section{Estratégias de escalabilidade}

\subsection{Automação e tecnologia}

Escalar significa atender mais tutores sem aumentar os custos na mesma proporção. O primeiro passo é padronizar: ficha de parceiro, descrição de serviços, horários, critérios de preço social e fluxo de encaminhamento. Depois, automatizar apenas o que se repete e já é compreendido, como lembretes consentidos, confirmação de disponibilidade e relatórios para parceiros. Mapa real, agenda e intermediação financeira dependem de integração e testes; a tela atual não comprova essas capacidades.

\subsection{Modelo replicável por território}

O modelo é replicável como plataforma digital: cada novo território reaproveita o mesmo aplicativo, os mesmos processos e o mesmo material de integração de parceiros, com custo adicional menor [10]. A regra de entrada é ter oferta mínima confirmada e economia unitária compreendida. Hoffman e Yeh descrevem a escala acelerada como estratégia para mercados em que o vencedor leva quase tudo [27]; esse não é o caso do PetCuida nesta fase, e franquias ou licenciamento ficam para depois de um processo estável e de revisão jurídica.

\subsection{Receita recorrente}

As fontes recorrentes propostas são o plano mensal para parceiros e convênios ou patrocínios com prazo definido. Em plataformas de dois lados, é comum cobrar do lado que obtém mais valor e subsidiar o lado mais sensível a preço [9]; por isso, as funções essenciais de cuidado preventivo permanecem gratuitas para o tutor.

\subsection{Dados e segurança}

Antes do piloto, a equipe deve definir finalidade e consentimento, limitar a coleta, controlar acessos e prever resposta a incidentes, conforme a LGPD e o guia da ANPD para agentes de pequeno porte [26, 28].
